\chapter{Solution Concept}
\label{chap:concept}

On the basis of the problem analysis of Chapter~3 and the theoretical
knowledge of Chapter~2, this chapter develops the solution concept.
Section~4.1 gives an overview of the system architecture. Section~4.2
specifies the Mixture of Experts model, Section~4.3 the router, and
Section~4.4 the continual training procedure.

\section{Architectural Overview}
\label{sec:overview}

The system is organised in four loosely coupled layers, which are shown in
Figure~4.1 together with their data flow.

\begin{figure}[H]
\centering
\begin{tikzpicture}[
  box/.style={draw, rounded corners, minimum height=1.05cm, minimum width=2.75cm, align=center, fill=blue!6, font=\small},
  arr/.style={-{Stealth[length=2.5mm]}, thick}]
  \node[box] (data) {Benchmark layer\\[1pt] \scriptsize SplitMNIST,\\ \scriptsize PermutedMNIST};
  \node[box, right=0.6cm of data] (model) {Model layer\\[1pt] \scriptsize MLP, Wide MLP, MoE};
  \node[box, right=0.6cm of model] (train) {Strategy layer\\[1pt] \scriptsize Naive, EWC, MoE};
  \node[box, right=0.6cm of train] (eval) {Evaluation layer\\[1pt] \scriptsize $R$ matrix, metrics};
  \draw[arr] (data) -- (model);
  \draw[arr] (model) -- (train);
  \draw[arr] (train) -- (eval);
\end{tikzpicture}
\caption{Layered concept of the experimental system. The data flows from
the benchmark layer through the model and strategy layers to the evaluation
layer, which produces the accuracy matrix, the scalar metrics and the
router profiles.}
\label{fig:system}
\end{figure}

The benchmark layer turns a static data set into a sequence of tasks and
provides one training and one test data loader per task. The model layer
contains the network architectures: the baseline MLP, the wide MLP and the
Mixture of Experts network. The strategy layer implements the training
procedures, namely naive fine tuning, EWC and the MoE training with the
auxiliary loss. Finally, the evaluation layer computes the accuracy matrix
$R$ after every training stage, derives the scalar metrics from it and
extracts the router profiles for the MoE model.

This separation had two practical advantages during the project. First,
every component could be implemented and tested independently. Second, the
comparison between the approaches stays fair, because all approaches
consume exactly the same task sequence and are evaluated by exactly the
same code. This directly answers the fairness sub problem SP4 of
Section~3.1.

\section{Concept of the MoE Model}
\label{sec:moedesign}

The proposed model follows the structure sketched in Figure~4.2. A shared
trunk maps the flattened $28 \times 28$ input image to a $400$ dimensional
representation. This representation is processed by one MoE layer
consisting of $E = 8$ experts and a router. Each expert is a small feed
forward block with one hidden layer and ReLU activation, with identical
input and output dimension, so that the expert outputs can be combined by a
weighted sum. The router selects the $k = 2$ experts with the highest gate
values for every sample and combines their outputs according to
Equation~\eqref{eq:moe}. On top of the MoE layer, one linear output head
per task produces the class logits (multi head protocol, see
Section~2.2.2).

\begin{figure}[H]
\centering
\begin{tikzpicture}[
  x=1cm, y=1cm,
  io/.style={draw, minimum width=1.15cm, minimum height=1.7cm, fill=gray!10, align=center},
  exp/.style={draw, minimum width=0.75cm, minimum height=1.25cm, fill=blue!12},
  expa/.style={draw, minimum width=0.75cm, minimum height=1.25cm, fill=orange!30},
  router/.style={draw, trapezium, trapezium angle=70, minimum width=1.9cm, fill=green!15, align=center},
  sum/.style={draw, circle, inner sep=1.5pt, fill=red!15},
  arr/.style={-{Stealth[length=2.2mm]}}]
  \node[io] (in) at (0,0) {input\\ $784$};
  \node[io, fill=gray!20] (trunk) at (2.2,0) {trunk\\ $784{\rightarrow}400$};
  \node[router] (gate) at (4.6,1.9) {\footnotesize router};
  \node[exp]  (e1) at (4.6,0.9) {};
  \node[expa] (e2) at (4.6,0.15) {};
  \node at (4.6,-0.5) {\footnotesize $\vdots$};
  \node[expa] (e8) at (4.6,-1.1) {};
  \node at (5.6,0) {\footnotesize $E=8$ experts};
  \node[sum] (plus) at (7.1,0) {$\sum$};
  \node[io, fill=gray!20] (head) at (9.2,0) {task head\\ $400{\rightarrow}c$};
  \node[io] (out) at (11.2,0) {logits};
  \draw[arr] (in) -- (trunk);
  \draw[arr] (trunk) -- (gate);
  \draw[arr] (trunk) -- (4.0,0) -- (e1);
  \draw[arr] (4.0,0) -- (e2);
  \draw[arr] (4.0,0) -- (e8);
  \draw[arr,dashed] (gate) -- (e1);
  \draw[arr,dashed] (gate) -- (e2);
  \draw[arr,dashed] (gate) -- (e8);
  \draw[arr] (e1) -- (plus);
  \draw[arr] (e2) -- (plus);
  \draw[arr] (e8) -- (plus);
  \draw[arr] (plus) -- (head);
  \draw[arr] (head) -- (out);
\end{tikzpicture}
\caption{Concept of the MoE model. A shared trunk feeds eight parallel
experts. The router (dashed arrows) selects the top 2 experts per sample
(highlighted); their outputs are combined as a gate weighted sum and passed
to the output head of the current task.}
\label{fig:moearch}
\end{figure}

The choice $E = 8$ with $k = 2$ follows the efficiency requirement of
Section~1.2: with five tasks, eight experts give enough spare capacity for
the router to separate the tasks, while the activated computation per
sample, namely two experts, stays close to that of the baseline. The output
head per task keeps the comparison with the baselines protocol fair,
because the baselines use identical heads.

\section{Concept of the Router}
\label{sec:routerdesign}

The router is a single linear layer $W_g \in \mathbb{R}^{8 \times 400}$
followed by a softmax, see Equation~\eqref{eq:gate}. It is deliberately
kept simple: a small gate cannot itself become a bottleneck of forgetting,
and its decisions stay easy to interpret and to visualise. Two design
decisions are essential for the continual setting:

\begin{enumerate}[label=\textbf{D\arabic*.}, leftmargin=3em]
  \item The router does not receive the task identity. The gate only sees
        the hidden representation of the input. Every task separation it
        shows is therefore learned from the data alone, which makes the
        router analysis in Chapter~7 meaningful.
  \item The router is trained only on the current task, without any
        mechanism that forces it to keep old routing decisions. If the
        gate changed its assignments freely, previously acquired expert
        knowledge would be re-routed and lost. The experiments test whether
        the assignments stay stable enough in practice.
\end{enumerate}

The auxiliary loss for load balancing from Equation~\eqref{eq:lb}, with the
coefficient $\alpha = 0.01$, prevents expert collapse at the beginning of
training and guarantees that all eight experts stay functional, so that
later tasks can actually find unused capacity. As the evaluation will show,
this component behaves differently than expected in the continual setting,
which leads to the architecture variant described in Section~7.5.

\paragraph{Considered alternatives.} Two alternative router designs were
considered and rejected. The first alternative was a nonlinear gate with
one hidden layer. It was rejected because the gate itself is part of the
continual learning problem: a larger gate has more parameters that can
forget, and its decisions are harder to interpret in the router analysis.
The second alternative was a hard top one routing without any mixing,
as in the Switch Transformer \cite{fedus2022}. It was rejected for this
small scale setting because the gradient signal of a single expert per
sample is noisy, and because a soft mixing of two experts makes the
transition between tasks smoother when the router changes its assignment.
The top two choice keeps most of the sparsity but keeps training stable,
which was confirmed in preliminary tests.

\paragraph{Expected behaviour.} If the concept works as intended, the
following behaviour should be observable in the experiments. During the
first task, the router distributes the samples over several experts. When
the second task arrives, the router should prefer experts that carry little
knowledge of the first task, so that the gradient updates of the new task
mainly rewrite unused or weakly used parameters. Over the whole sequence,
a pattern of task specific expert subsets should emerge, and the pairwise
routing similarity between tasks should stay low. These predictions are
checked explicitly in Section~7.6.

\section{Concept of the Training Procedure}
\label{sec:trainingdesign}

All approaches follow the same sequential protocol, which is shown in
Table~4.1. For every task, the model is trained from its current state on
the training data of that task; afterwards, it is evaluated on the test
data of all tasks seen so far, which produces one row of the accuracy
matrix $R$. No data of earlier tasks is shown again, and no task boundaries
are signalled to the router.

\begin{table}[H]
\centering
\caption{Sequential continual training protocol used for all approaches.}
\label{tab:protocol}
\begin{tabular}{p{0.86\textwidth}}
\toprule
\textbf{Algorithm:} Sequential training over a task sequence \\
\midrule
\textbf{Input:} task sequence $(\mathcal{T}_1, \dots, \mathcal{T}_T)$;
model $f_\theta$; strategy specific loss $\mathcal{L}$ \\
1: \textbf{for} $t = 1$ \textbf{to} $T$ \textbf{do} \\
2: \quad train $f_\theta$ on $\mathcal{T}_t$ by minimising $\mathcal{L}$ \\
3: \quad \textbf{for} $j = 1$ \textbf{to} $t$ \textbf{do} \\
4: \quad \quad $R_{t,j} \leftarrow$ accuracy of $f_\theta$ on the test set of $\mathcal{T}_j$ \\
5: \quad \textbf{end for} \\
6: \textbf{end for} \\
7: derive $A$, $F$ and BWT from $R$ (Equations \eqref{eq:avgacc} to \eqref{eq:bwt}) \\
\bottomrule
\end{tabular}
\end{table}

For the EWC strategy, step 2 additionally estimates the diagonal Fisher
information on the current task after convergence and accumulates the
quadratic penalty of Equation~\eqref{eq:ewc} for all later tasks. For the
MoE strategy, step 2 minimises the combined loss of
Equation~\eqref{eq:totalloss}. The evaluation in lines 3 to 5 is identical
for all approaches, which guarantees that differences in the metrics can be
attributed to the models and strategies alone.
